% --- Archon physics formalization source begin ---
% archon:physics
% archon:covers PhyXMiniProblems/problem_phyx_mini_0384.lean
% archon:source-report reports/phyx_mini/problem_phyx_mini_0384.source.json

\chapter{Physics problem phyx\_mini\_0384}
\label{ch:PhyXMiniProblems_problem_phyx_mini_0384}

\paragraph{Problem source.}
\#\# Physical scenario

A piston/cylinder with a cross-sectional area of 0.01 \textbackslash{}m\textasciicircum{}2\textbackslash{} has a piston mass of 100 kg resting on the stops, as shown in the figure.

\#\# Question

With an outside atmospheric pressure of 100 kPa, what should the water pressure be to lift the piston?

\#\# Answer choices

- A: A: 490kPa

- B: B: 198kPa

- C: C: 154kPa

- D: D: 10.2kPa

\#\# Figure caption (auxiliary; use the image as primary evidence)

The image is a diagram of a container filled with water. Above the water, there's a horizontal piston inside the container. The piston appears to be sealing the container at the top, and there are downward arrows indicating the force of gravity labeled as "g."

Text elements in the image include:

- "Water" is written inside the container, denoting the liquid's content.
- "P₀" is written above the piston, indicating the atmospheric or external pressure exerted on the piston.

The components are enclosed within walls, and the setup suggests a study of fluid mechanics involving pressure and gravity interactions with the water in the container.

\#\# Dataset metadata

Laws of Thermodynamics; Physical Model Grounding Reasoning, Multi-Formula Reasoning

\paragraph{Recorded answer/context.}
B: B: 198kPa

\paragraph{Figure/image path.}
/root/proposal\_for\_physic/science-mango/phyx\_mini\_run/phyx\_data/test\_image/384.png

\paragraph{Formalization target.}
create a compiling Lean file with sorry bodies at `PhyXMiniProblems/problem\_phyx\_mini\_0384.lean`. The Lean declarations must preserve the physical quantities, dimensions or dimensional roles, figure labels, governing-law hypotheses, and final relation expressed by this problem.
Use Mathlib/Physlib names found through LeanExplore where available. If a physics API is missing, introduce faithful local abstractions rather than scalar placeholder aliases.

\begin{theorem}[Physics formalization target]
\label{thm:physics:phyx_mini_0384:target}
\lean{PhyXMiniProblems.ProblemPhyXMini0384.waterPressureToLiftPiston_eq_198_kPa}
\uses{def:physics:phyx-mini-0384:phyxminiproblems-problemphyxmini0384-pressureinkilopascals, def:physics:phyx-mini-0384:phyxminiproblems-problemphyxmini0384-pistoncylindersetup, def:physics:phyx-mini-0384:phyxminiproblems-problemphyxmini0384-matchesprimaryfigure, def:physics:phyx-mini-0384:phyxminiproblems-problemphyxmini0384-matchesproblemdata, def:physics:phyx-mini-0384:phyxminiproblems-problemphyxmini0384-haspositivephysicalparameters, def:physics:phyx-mini-0384:phyxminiproblems-problemphyxmini0384-atincipientlift}
The assigned autoformalize agent should translate this physics problem into Lean declarations in the covered file, with theorem and lemma proof bodies written as `by sorry`.
\end{theorem}
\begin{proof}
This is an autoformalization task, not a proof task. Produce statements that can later be proved by the physics prover without weakening the physical model.
\end{proof}

% --- Archon declaration topology begin ---
\section{Declaration topology}
\begin{definition}[Mass Quantity]
  \label{def:physics:phyx-mini-0384:phyxminiproblems-problemphyxmini0384-massquantity}
  \lean{PhyXMiniProblems.ProblemPhyXMini0384.MassQuantity}
  A nonnegative physical mass, independent of the unit used to read it
\end{definition}
\begin{proof}
  This is the declared model definition of Mass Quantity.
\end{proof}
\begin{definition}[Acceleration Quantity]
  \label{def:physics:phyx-mini-0384:phyxminiproblems-problemphyxmini0384-accelerationquantity}
  \lean{PhyXMiniProblems.ProblemPhyXMini0384.AccelerationQuantity}
  A nonnegative physical acceleration magnitude
\end{definition}
\begin{proof}
  This is the declared model definition of Acceleration Quantity.
\end{proof}
\begin{definition}[Area Quantity]
  \label{def:physics:phyx-mini-0384:phyxminiproblems-problemphyxmini0384-areaquantity}
  \lean{PhyXMiniProblems.ProblemPhyXMini0384.AreaQuantity}
  A nonnegative physical cross-sectional area
\end{definition}
\begin{proof}
  This is the declared model definition of Area Quantity.
\end{proof}
\begin{definition}[Pressure Quantity]
  \label{def:physics:phyx-mini-0384:phyxminiproblems-problemphyxmini0384-pressurequantity}
  \lean{PhyXMiniProblems.ProblemPhyXMini0384.PressureQuantity}
  A physical absolute pressure
\end{definition}
\begin{proof}
  This is the declared model definition of Pressure Quantity.
\end{proof}
\begin{definition}[mass In Kilograms]
  \label{def:physics:phyx-mini-0384:phyxminiproblems-problemphyxmini0384-massinkilograms}
  \lean{PhyXMiniProblems.ProblemPhyXMini0384.massInKilograms}
  \uses{def:physics:phyx-mini-0384:phyxminiproblems-problemphyxmini0384-massquantity}
  Coherent-SI kilogram readout of a physical mass
\end{definition}
\begin{proof}
  This is the declared model definition of mass In Kilograms.
\end{proof}
\begin{definition}[acceleration In Meters Per Second Squared]
  \label{def:physics:phyx-mini-0384:phyxminiproblems-problemphyxmini0384-accelerationinmeterspersecondsquared}
  \lean{PhyXMiniProblems.ProblemPhyXMini0384.accelerationInMetersPerSecondSquared}
  \uses{def:physics:phyx-mini-0384:phyxminiproblems-problemphyxmini0384-accelerationquantity}
  Coherent-SI metre-per-second-squared readout of an acceleration
\end{definition}
\begin{proof}
  This is the declared model definition of acceleration In Meters Per Second Squared.
\end{proof}
\begin{definition}[area In Square Meters]
  \label{def:physics:phyx-mini-0384:phyxminiproblems-problemphyxmini0384-areainsquaremeters}
  \lean{PhyXMiniProblems.ProblemPhyXMini0384.areaInSquareMeters}
  \uses{def:physics:phyx-mini-0384:phyxminiproblems-problemphyxmini0384-areaquantity}
  Coherent-SI square-metre readout of an area
\end{definition}
\begin{proof}
  This is the declared model definition of area In Square Meters.
\end{proof}
\begin{definition}[pressure In Pascals]
  \label{def:physics:phyx-mini-0384:phyxminiproblems-problemphyxmini0384-pressureinpascals}
  \lean{PhyXMiniProblems.ProblemPhyXMini0384.pressureInPascals}
  \uses{def:physics:phyx-mini-0384:phyxminiproblems-problemphyxmini0384-pressurequantity}
  Coherent-SI pascal readout of an absolute pressure
\end{definition}
\begin{proof}
  This is the declared model definition of pressure In Pascals.
\end{proof}
\begin{definition}[pressure In Kilopascals]
  \label{def:physics:phyx-mini-0384:phyxminiproblems-problemphyxmini0384-pressureinkilopascals}
  \lean{PhyXMiniProblems.ProblemPhyXMini0384.pressureInKilopascals}
  \uses{def:physics:phyx-mini-0384:phyxminiproblems-problemphyxmini0384-pressurequantity, def:physics:phyx-mini-0384:phyxminiproblems-problemphyxmini0384-pressureinpascals}
  Kilopascal readout used by the four displayed answers
\end{definition}
\begin{proof}
  This is the declared model definition of pressure In Kilopascals.
\end{proof}
\begin{definition}[Cylinder Fluid]
  \label{def:physics:phyx-mini-0384:phyxminiproblems-problemphyxmini0384-cylinderfluid}
  \lean{PhyXMiniProblems.ProblemPhyXMini0384.CylinderFluid}
  The fluid occupying the cylinder below the piston
\end{definition}
\begin{proof}
  This is the declared model definition of Cylinder Fluid.
\end{proof}
\begin{definition}[Figure Region]
  \label{def:physics:phyx-mini-0384:phyxminiproblems-problemphyxmini0384-figureregion}
  \lean{PhyXMiniProblems.ProblemPhyXMini0384.FigureRegion}
  Regions distinguished by the supplied cross-sectional diagram
\end{definition}
\begin{proof}
  This is the declared model definition of Figure Region.
\end{proof}
\begin{definition}[Vertical Direction]
  \label{def:physics:phyx-mini-0384:phyxminiproblems-problemphyxmini0384-verticaldirection}
  \lean{PhyXMiniProblems.ProblemPhyXMini0384.VerticalDirection}
  Vertical direction of a force or acceleration arrow
\end{definition}
\begin{proof}
  This is the declared model definition of Vertical Direction.
\end{proof}
\begin{definition}[Initial Piston Support]
  \label{def:physics:phyx-mini-0384:phyxminiproblems-problemphyxmini0384-initialpistonsupport}
  \lean{PhyXMiniProblems.ProblemPhyXMini0384.InitialPistonSupport}
  The piston's support condition before the water pressure is increased
\end{definition}
\begin{proof}
  This is the declared model definition of Initial Piston Support.
\end{proof}
\begin{definition}[Piston Cylinder Figure]
  \label{def:physics:phyx-mini-0384:phyxminiproblems-problemphyxmini0384-pistoncylinderfigure}
  \lean{PhyXMiniProblems.ProblemPhyXMini0384.PistonCylinderFigure}
  \uses{def:physics:phyx-mini-0384:phyxminiproblems-problemphyxmini0384-figureregion, def:physics:phyx-mini-0384:phyxminiproblems-problemphyxmini0384-verticaldirection}
  Labels and placements read directly from the primary figure
\end{definition}
\begin{proof}
  This is the declared model definition of Piston Cylinder Figure.
\end{proof}
\begin{definition}[Piston Cylinder Setup]
  \label{def:physics:phyx-mini-0384:phyxminiproblems-problemphyxmini0384-pistoncylindersetup}
  \lean{PhyXMiniProblems.ProblemPhyXMini0384.PistonCylinderSetup}
  \uses{def:physics:phyx-mini-0384:phyxminiproblems-problemphyxmini0384-massquantity, def:physics:phyx-mini-0384:phyxminiproblems-problemphyxmini0384-accelerationquantity, def:physics:phyx-mini-0384:phyxminiproblems-problemphyxmini0384-areaquantity, def:physics:phyx-mini-0384:phyxminiproblems-problemphyxmini0384-pressurequantity, def:physics:phyx-mini-0384:phyxminiproblems-problemphyxmini0384-cylinderfluid, def:physics:phyx-mini-0384:phyxminiproblems-problemphyxmini0384-initialpistonsupport, def:physics:phyx-mini-0384:phyxminiproblems-problemphyxmini0384-pistoncylinderfigure}
  The piston/cylinder and the physical quantities relevant at the threshold of lift. liftThresholdWaterPressure is an unknown physical pressure field; no definition assigns it the requested numerical value
\end{definition}
\begin{proof}
  This is the declared model definition of Piston Cylinder Setup.
\end{proof}
\begin{definition}[Matches Primary Figure]
  \label{def:physics:phyx-mini-0384:phyxminiproblems-problemphyxmini0384-matchesprimaryfigure}
  \lean{PhyXMiniProblems.ProblemPhyXMini0384.MatchesPrimaryFigure}
  \uses{def:physics:phyx-mini-0384:phyxminiproblems-problemphyxmini0384-pistoncylindersetup}
  The apparatus and labels visible in the source diagram
\end{definition}
\begin{proof}
  This is the declared model definition of Matches Primary Figure.
\end{proof}
\begin{definition}[Matches Problem Data]
  \label{def:physics:phyx-mini-0384:phyxminiproblems-problemphyxmini0384-matchesproblemdata}
  \lean{PhyXMiniProblems.ProblemPhyXMini0384.MatchesProblemData}
  \uses{def:physics:phyx-mini-0384:phyxminiproblems-problemphyxmini0384-massinkilograms, def:physics:phyx-mini-0384:phyxminiproblems-problemphyxmini0384-accelerationinmeterspersecondsquared, def:physics:phyx-mini-0384:phyxminiproblems-problemphyxmini0384-areainsquaremeters, def:physics:phyx-mini-0384:phyxminiproblems-problemphyxmini0384-pressureinpascals, def:physics:phyx-mini-0384:phyxminiproblems-problemphyxmini0384-pistoncylindersetup}
  Numerical readouts supplied by the prose, together with the standard g = 9.8 m/s² value used by the recorded multiple-choice calculation
\end{definition}
\begin{proof}
  This is the declared model definition of Matches Problem Data.
\end{proof}
\begin{definition}[Has Positive Physical Parameters]
  \label{def:physics:phyx-mini-0384:phyxminiproblems-problemphyxmini0384-haspositivephysicalparameters}
  \lean{PhyXMiniProblems.ProblemPhyXMini0384.HasPositivePhysicalParameters}
  \uses{def:physics:phyx-mini-0384:phyxminiproblems-problemphyxmini0384-massinkilograms, def:physics:phyx-mini-0384:phyxminiproblems-problemphyxmini0384-accelerationinmeterspersecondsquared, def:physics:phyx-mini-0384:phyxminiproblems-problemphyxmini0384-areainsquaremeters, def:physics:phyx-mini-0384:phyxminiproblems-problemphyxmini0384-pressureinpascals, def:physics:phyx-mini-0384:phyxminiproblems-problemphyxmini0384-pistoncylindersetup}
  Positivity assumptions selecting the intended physical branch
\end{definition}
\begin{proof}
  This is the declared model definition of Has Positive Physical Parameters.
\end{proof}
\begin{definition}[upward Water Pressure Force In Newtons]
  \label{def:physics:phyx-mini-0384:phyxminiproblems-problemphyxmini0384-upwardwaterpressureforceinnewtons}
  \lean{PhyXMiniProblems.ProblemPhyXMini0384.upwardWaterPressureForceInNewtons}
  \uses{def:physics:phyx-mini-0384:phyxminiproblems-problemphyxmini0384-areainsquaremeters, def:physics:phyx-mini-0384:phyxminiproblems-problemphyxmini0384-pressureinpascals, def:physics:phyx-mini-0384:phyxminiproblems-problemphyxmini0384-pistoncylindersetup}
  Upward force readout, in newtons, exerted by the water on the piston
\end{definition}
\begin{proof}
  This is the declared model definition of upward Water Pressure Force In Newtons.
\end{proof}
\begin{definition}[downward Atmospheric Force In Newtons]
  \label{def:physics:phyx-mini-0384:phyxminiproblems-problemphyxmini0384-downwardatmosphericforceinnewtons}
  \lean{PhyXMiniProblems.ProblemPhyXMini0384.downwardAtmosphericForceInNewtons}
  \uses{def:physics:phyx-mini-0384:phyxminiproblems-problemphyxmini0384-areainsquaremeters, def:physics:phyx-mini-0384:phyxminiproblems-problemphyxmini0384-pressureinpascals, def:physics:phyx-mini-0384:phyxminiproblems-problemphyxmini0384-pistoncylindersetup}
  Downward force readout, in newtons, due to the outside atmosphere
\end{definition}
\begin{proof}
  This is the declared model definition of downward Atmospheric Force In Newtons.
\end{proof}
\begin{definition}[piston Weight In Newtons]
  \label{def:physics:phyx-mini-0384:phyxminiproblems-problemphyxmini0384-pistonweightinnewtons}
  \lean{PhyXMiniProblems.ProblemPhyXMini0384.pistonWeightInNewtons}
  \uses{def:physics:phyx-mini-0384:phyxminiproblems-problemphyxmini0384-massinkilograms, def:physics:phyx-mini-0384:phyxminiproblems-problemphyxmini0384-accelerationinmeterspersecondsquared, def:physics:phyx-mini-0384:phyxminiproblems-problemphyxmini0384-pistoncylindersetup}
  Downward gravitational force readout, in newtons, on the piston
\end{definition}
\begin{proof}
  This is the declared model definition of piston Weight In Newtons.
\end{proof}
\begin{definition}[At Incipient Lift]
  \label{def:physics:phyx-mini-0384:phyxminiproblems-problemphyxmini0384-atincipientlift}
  \lean{PhyXMiniProblems.ProblemPhyXMini0384.AtIncipientLift}
  \uses{def:physics:phyx-mini-0384:phyxminiproblems-problemphyxmini0384-pistoncylindersetup, def:physics:phyx-mini-0384:phyxminiproblems-problemphyxmini0384-upwardwaterpressureforceinnewtons, def:physics:phyx-mini-0384:phyxminiproblems-problemphyxmini0384-downwardatmosphericforceinnewtons, def:physics:phyx-mini-0384:phyxminiproblems-problemphyxmini0384-pistonweightinnewtons}
  The ideal vertical force balance at the instant the piston loses contact with the stops. This governing law includes no numerical value for the unknown water pressure and no answer-choice label
\end{definition}
\begin{proof}
  This is the declared model definition of At Incipient Lift.
\end{proof}
\begin{definition}[Answer Choice]
  \label{def:physics:phyx-mini-0384:phyxminiproblems-problemphyxmini0384-answerchoice}
  \lean{PhyXMiniProblems.ProblemPhyXMini0384.AnswerChoice}
  Labels of the four answers printed with the problem
\end{definition}
\begin{proof}
  This is the declared model definition of Answer Choice.
\end{proof}
\begin{definition}[answer Pressure In Kilopascals]
  \label{def:physics:phyx-mini-0384:phyxminiproblems-problemphyxmini0384-answerpressureinkilopascals}
  \lean{PhyXMiniProblems.ProblemPhyXMini0384.answerPressureInKilopascals}
  \uses{def:physics:phyx-mini-0384:phyxminiproblems-problemphyxmini0384-answerchoice}
  Kilopascal value printed beside each answer label
\end{definition}
\begin{proof}
  This is the declared model definition of answer Pressure In Kilopascals.
\end{proof}
\begin{definition}[recorded Answer Choice]
  \label{def:physics:phyx-mini-0384:phyxminiproblems-problemphyxmini0384-recordedanswerchoice}
  \lean{PhyXMiniProblems.ProblemPhyXMini0384.recordedAnswerChoice}
  \uses{def:physics:phyx-mini-0384:phyxminiproblems-problemphyxmini0384-answerchoice}
  The answer label recorded in the supplied dataset
\end{definition}
\begin{proof}
  This is the declared model definition of recorded Answer Choice.
\end{proof}
% --- Archon declaration topology end ---
% --- Archon physics formalization source end ---
